% Parte 2 - Tabla de servicios y requisitos (SSDC). Fragmento: se incluye con \input desde main.tex.
% Requiere: booktabs, longtable, array y url (\usepackage{url}) para que las rutas largas se partan en la celda.
% Nombres de servicio segun la convencion ms-<dominio> de la guia del equipo
% (en el repositorio de E1 se llaman convocatorias, postulaciones, evaluacion, notificaciones, auditoria, ia).
% Marca (impl.) = endpoint ya implementado en el esqueleto de src/. Lo demas es PROPUESTA de diseno.
{\footnotesize
\begin{longtable}{@{}p{3.2cm}p{2.6cm}p{2.7cm}p{4.8cm}p{1.7cm}@{}}
\caption{Tabla de servicios del SSDC y requisitos funcionales que atienden}\\
\toprule
\textbf{Servicio} & \textbf{Responsabilidad} & \textbf{Datos que le pertenecen} & \textbf{Endpoints o eventos principales} & \textbf{RF que atiende} \\
\midrule
\endfirsthead
\toprule
\textbf{Servicio} & \textbf{Responsabilidad} & \textbf{Datos que le pertenecen} & \textbf{Endpoints o eventos principales} & \textbf{RF que atiende} \\
\midrule
\endhead
\path|ms-convocatorias| &
Administrar las convocatorias y el tr\'amite de aprobaci\'on del informe ante los consejos. &
Convocatorias, su estado y las instancias de aprobaci\'on. Base: \path|db-convocatorias|. &
\path|GET /api/v1/convocatorias| (impl.) \newline
\path|POST /api/v1/convocatorias/:id/publicar| (impl.) \newline
\path|POST /api/v1/convocatorias/:id/tramite| \newline
Publica \path|convocatorias.convocatoria_publicada.v1| &
RF-10 \\
\midrule
\path|ms-postulaciones| &
Recibir y guardar las postulaciones con su hoja de vida y mostrar su estado. &
Postulantes, expediente en PDF (referencia a Storage) y estado de la postulaci\'on. Base: \path|db-postulaciones|. &
\path|POST /api/v1/postulaciones| \newline
\path|GET /api/v1/convocatorias/:id/postulaciones| \newline
\path|GET /api/v1/postulaciones/:id/estado| \newline
Publica \path|postulaciones.postulacion_recibida.v1| y \path|postulaciones.estado_cambiado.v1| &
RF-01, RF-02, RF-03 \\
\midrule
\path|ms-evaluacion| &
Evaluar a los postulantes: requisitos m\'inimos, matriz de m\'eritos, entrevista, ranking e informe. &
Matriz y puntajes por criterio, entrevistas, total, ranking y marca de cumplimiento. Base: \path|db-evaluacion|. &
\path|PUT /api/v1/evaluaciones/:id/criterios| \newline
\path|GET /api/v1/evaluaciones/:id/total| \newline
\path|POST /api/v1/entrevistas| \newline
\path|GET /api/v1/convocatorias/:id/ranking| \newline
\path|POST /api/v1/convocatorias/:id/informe| \newline
Publica \path|evaluacion.evaluacion_terminada.v1|; consulta a \path|ms-postulaciones| &
RF-04, RF-05, RF-06, RF-07, RF-08, RF-09 \\
\midrule
\path|ms-notificaciones| &
Avisar a los usuarios los cambios de estado y las citaciones. &
Avisos generados y su estado de env\'io (enviado o pendiente). Base: \path|db-notificaciones|. &
\path|GET /api/v1/notificaciones| (impl. como \path|GET /notificaciones|) \newline
\path|POST /api/v1/notificaciones| (impl. como \path|POST /notificaciones|) \newline
Consume los eventos de postulaciones, evaluaci\'on y convocatorias &
RF-03, RF-07 (aviso) \\
\midrule
\path|ms-auditoria| &
Registrar los cambios de puntaje y de estado, y respaldarlos. &
Registros de cambio con hash y respaldos. Base: \path|db-auditoria|. &
\path|GET /api/v1/auditoria| \newline
Consume todos los eventos del sistema &
RF-12 (nuevo); tambi\'en RNF-04 \\
\midrule
\path|ms-ia| &
Generar textos e im\'agenes asistidos con proveedores externos de IA. &
Solicitudes de generaci\'on y su resultado. Base: \path|db-ia|. &
\path|POST /api/v1/ia/generaciones| &
RF-13 (nuevo) \\
\midrule
\path|api-gateway| (borde) &
Puerta \'unica de entrada: enrutar, autenticar (valida el token y el rol que trae) y limitar solicitudes. &
Ninguno (sin estado de negocio). &
Rutas \path|/api/v1/*| hacia cada servicio &
RF-11 (acceso por rol) \\
\bottomrule
\end{longtable}
}

\noindent\textbf{Requisitos nuevos acordados por el equipo} (redacci\'on aprobada por el equipo):
\begin{itemize}
\item \textbf{RF-12.} Como \textbf{administrador}, quiero consultar el registro de cambios de puntajes y estados (qui\'en, cu\'ando, valor anterior y nuevo), para verificar que las evaluaciones no fueron alteradas.
\item \textbf{RF-13.} Como \textbf{Coordinadora}, quiero generar con ayuda de IA el borrador del texto y la imagen de difusi\'on de una convocatoria, para no redactarlos desde cero.
\end{itemize}

\noindent\textbf{Verificaci\'on de cobertura (contra RF-01 a RF-13).}
\begin{itemize}
\item Requisitos sin servicio: ninguno. RF-04, RF-09 y RF-10 se asignaron por acuerdo del equipo; RF-12 y RF-13 son nuevos y los atienden \path|ms-auditoria| y \path|ms-ia|; RF-11 lo atiende el \path|api-gateway|.
\item Servicios sin requisito funcional: ninguno.
\end{itemize}
